
\section{Excel Shrink}
\label{sec:methodology4}


%\begin{algorithm}[ht]
%  \caption{Clean Excel Files with Multithreading}
%  \label{alg:multithreading}
%  \begin{algorithmic}
%
%    \STATE \textbf{Input:} 
%    \STATE \quad \texttt{InP} $\leftarrow$ Excel file path
%    \STATE \quad \texttt{OutP} $\leftarrow$ Output path
%    \vspace{0.2cm}
%    
%    \STATE \texttt{InArc} $\leftarrow$ Open the ZIP archive \texttt{InP} for reading
%    \STATE \texttt{OutArc} $\leftarrow$ Create a new ZIP archive for writing
%    \STATE \texttt{ShStr} $\leftarrow$ Extract shared strings from \texttt{sharedStrings.xml}
%    \STATE \texttt{Sty} $\leftarrow$ Extract styles from \texttt{styles.xml}
%    \STATE \texttt{SF} $\leftarrow$ List of all sheet XML files in \texttt{InArc}
%    \STATE \texttt{NSF} $\leftarrow$ List of non-sheet files in \texttt{InArc}    
%    \vspace{0.2cm}
%
%    \STATE \texttt{wvp} $\leftarrow$ dynamically create a pattern based on \texttt{ShStr} to identify whitespace-only cell values
%    \STATE \texttt{icp} $\leftarrow$ dynamically create a pattern based on \texttt{Sty} to identify cells that are invisible
%    \STATE \texttt{sdcp} $\leftarrow$ dynamically create a pattern based on \texttt{Sty} to identify cells whose styles depend on row/column styles 
%    \vspace{0.2cm}
%    \STATE \FOR{each sheet file \texttt{sf} in \texttt{SF} \textbf{in parallel}}
%        \STATE Clean sheet file \texttt{sf} in a stream-based manner using \texttt{wvp}, \texttt{icp}, \texttt{sdcp}, and \texttt{Sty}
%        \STATE Write cleaned sheet back to the ZIP archive \texttt{OutArc}
%    \ENDFOR
%
%    \STATE \FOR{each non-sheet file \texttt{nsf} in \texttt{NSF}}
%        \IF{\texttt{nsf} is \texttt{workbook.xml}}
%            \STATE Clean workbook metadata with a DOM-based approach
%            \STATE Write cleaned workbook to \texttt{OutArc}
%        \ELSE
%            \STATE Copy \texttt{nsf} unchanged to \texttt{OutArc}
%        \ENDIF
%    \ENDFOR
%
%    \STATE \textbf{Output:} Cleaned Excel file
%
%  \end{algorithmic}
%\end{algorithm}





% Our methodology evolved from an initial DOM-based prototype, the \emph{Excel Shrink Analyzer}, which helped us investigate and understand the primary causes of file bloat in large Excel files. This analysis revealed that excessive file sizes primarily resulted from hidden rows, empty rows, empty but styled cells, whitespace cells, and unnecessary columns.

% Because we encountered frequent memory limitations and poor scalability using XML-based approaches, we shifted to a \emph{stream-based approach}. This change was motivated by the insight that an efficient processing model could be implemented using \emph{regex-based chunk processing} on the raw bytestream.

% Constructing effective regular expressions required complex patterns, given the hierarchical nature of Excel’s XML structure. Consequently, we developed a pipeline of multiple regex stages, each designed to target a specific inefficiency in spreadsheet data:

% 1. Identify whitespace-only shared strings and remove their referencing cells.
% 2. Delete purely empty cells lacking any borders or fill styles.
% 3. Delete purely empty rows lacking any borders or fill styles nor a height.
% 4. Truncate rows that fall outside the bounding box of the visible data region.
% 5. Prune unnecessary column definitions, ensuring the final column range still extends to previous the maximum column index to preserve layout.

% This systematic approach enables efficient, scalable file-size reduction without requiring large in-memory representations, thereby significantly improving both performance and storage efficiency.





% \subsection{Parallel Context Extraction for Large-Scale Preprocessing}

% Both the pre-processing and post-processing phases are implemented using multithreading to improve throughput when handling large collections of Excel workbooks. This is only possible because we dont use xml parsing which would be memory consuming per se and would consume the whole systems memoery just by runing single threaded just as our xml streaming prof of concept did and because we dont use xml parsing but rahter perform operations on the raw bytestring and morover because we split huge document into smaller chunks only a negligible amount of memory is used per thread. The pipeline begins by parallelising the extraction of essential context information from each file. Given that Excel .xlsx files are essentially ZIP archives, they can be opened and inspected quickly using lightweight I/O operations.

\subsection{Multithreading for Performance}
To further increase parsing speed, we use multithreading. In this approach, each XML file within the ZIP archive is processed by its own thread, as described in Algorithm~\ref{alg:multithreading}. Because each thread splits the large \texttt{worksheetX.xml} files into smaller chunks and processes these chunks incrementally, only a negligible amount of memory is used per thread. This setup represents a significant improvement over our previous attepts using XML- DOM-based approach, where a single thread attempted to process all worksheets consecutively. The DOM-based method consistently consumed nearly 100\% of the available memory on our test machine, which prevented any realistic automation. By contrast, our new approach effectively reduces overall computation time and allows the files to be processed more quickly.

In this initial step, each thread is assigned one workbook and performs the following tasks:

Workbook Discovery: The internal structure of the archive is read to determine which worksheets are present.

Shared Strings Extraction: The sharedStrings.xml file is parsed to extract the mapping of shared string indices to their actual content. This information is later used to resolve cell values, particularly for detecting whitespace-only entries.

Style Information Parsing: The styles.xml file is analysed to extract all available styles, including fill and border definitions, as well as flags indicating whether these styles apply at the cell, row, or column level.

This extraction process is fully parallelised across all available CPU threads, with one thread per file. The number of concurrent threads is configured to match the number of virtual cores available on the machine, ensuring optimal utilisation of the host system’s hardware resources.

By performing context extraction up front in a multithreaded fashion, we ensure that all subsequent cleaning steps—many of which depend on shared strings and style definitions—can be executed without repeated I/O or redundant parsing effort.




%\begin{algorithm}[ht]
%  \caption{State-Machine Overview}
%  \label{alg:statemachine}
%  \begin{algorithmic}
%
%    \STATE \textbf{Input:} 
%    \STATE \quad \texttt{SheetS} $\leftarrow$ Sheet XML stream
%    \STATE \textbf{Initialize:}
%    \STATE \quad \texttt{Buffer} $\leftarrow$ Empty buffer
%    \STATE \quad \texttt{Output} $\leftarrow$ Empty buffer
%    \STATE \quad \texttt{State} $\leftarrow$ BEFORE\_COLS
%    \STATE \quad \texttt{MaxRow} $\leftarrow$ 0
%    \STATE \quad \texttt{MaxCol} $\leftarrow$ 0
%    \STATE \textbf{Processing Sheet Stream}
%    \WHILE{\texttt{State} $\neq$ END\_OF\_FILE}
%        \STATE Read chunk from \texttt{SheetS} into \texttt{Buffer}
%
%        \STATE \IF{\texttt{State} is in BEFORE\_COLS, BEFORE\_SHEET\_DATA \textbf{or} AFTER\_SHEET\_DATA} 
%            \IF{Found \texttt{<cols>} or \texttt{<sheetData>}}
%                \STATE Store \texttt{Buffer} in \texttt{Output} until last found tag
%                \STATE Discard content in \texttt{Buffer} until last found tag
%            \ENDIF
%            \IF{Found \texttt{<cols>}}
%              \STATE \texttt{State} $\leftarrow$ IN\_COLS
%            \ELSIF{Found \texttt{<sheetData>}}
%                \STATE \texttt{State} $\leftarrow$ IN\_SHEET\_DATA
%            \ELSIF{Reached end of file}
%                \STATE Store entire \texttt{Buffer} in \texttt{Output} 
%                \STATE \texttt{State} $\leftarrow$ END\_OF\_FILE
%            \ENDIF
%        \ELSIF{\texttt{State} is IN\_COLS}
%            \IF{Found \texttt{</cols>} or \texttt{<col/>}}
%                \STATE Parse column definitions
%                \STATE Store \texttt{Buffer} in \texttt{Output} until last found tag
%                \STATE Discard content in \texttt{Buffer} until last found tag
%            \ENDIF
%            \IF{Found \texttt{</cols>}}
%                \STATE \texttt{State} $\leftarrow$ BEFORE\_SHEET\_DATA
%            \ENDIF
%        \ELSIF{\texttt{State} is IN\_SHEET\_DATA}
%            \IF{Found \texttt{</row>}, \texttt{<row .../>} or \texttt{</sheetData>}}
%                \STATE Clean \texttt{Buffer} of invisible cells/rows using patterns
%                \STATE Store cleaned \texttt{Buffer} in \texttt{Output}
%                \STATE Discard content in \texttt{Buffer} until last found tag
%                \STATE Extract any remaining cell references to update \texttt{MaxRow} and \texttt{MaxCol}
%            \ENDIF
%            \IF{Found \texttt{</sheetData>}}
%                \STATE \texttt{State} $\leftarrow$ AFTER\_SHEET\_DATA
%            \ENDIF
%        \ENDIF
%    \ENDWHILE
%
%    \STATE Extract and update \texttt{MaxRow} and \texttt{MaxCol} from merge cells
%    \STATE Remove rows exceeding \texttt{MaxRow} from \texttt{sheetData}
%    \STATE Merge extraneous columns beyond \texttt{MaxCol}
%    \STATE \textbf{Output:} Cleaned sheet XML
%
%  \end{algorithmic}
%\end{algorithm}

\subsection{Methodology (High-Level State-Machine Overview)}


% As mentioned above, our new approach employs a state machine that helps manage the complexity of the operations needed to segment the file and apply different regular expressions.
% The state machine is split into multiple states, each representing a specific file section under processing, as described in Algorithm~\ref{alg:statemachine}, which provides a high-level overview of the state machine’s operation.
% The state machine transitions between states based on the XML content it encounters, allowing it to apply different regex patterns depending on the context. This approach ensures that the regex patterns are applied only to the relevant sections of the XML, preventing overmatching and preserving the XML structure.

%We implement our spreadsheet parser streaming-based, processing the XML data incrementally to minimize memory consumption and permit early discarding of unnecessary file sections. Both our “consecutive” and “interleaved” parsing approaches, which reduce memory usage, use this same underlying state machine. Below, we focus on how the parser transitions between states (i.e., how each state function decides to proceed), deferring lower-level details of in-state processing until later sections.
%
%\paragraph{States Overview.}
%We define six main states in our parser:
%\begin{itemize}
%  \item \texttt{BEFORE\_COLS}: The parser has not yet encountered the \texttt{<cols>} element. It searches for either \texttt{<cols>} or \texttt{<sheetData>}.
%  \item \texttt{IN\_COLS}: The parser found \texttt{<cols>} and is processing column definitions.
%  \item \texttt{BEFORE\_SHEET\_DATA}: The parser has finished reading (or did not detect) the \texttt{<cols>} node and is searching for the first \texttt{<sheetData>} element.
%  \item \texttt{IN\_SHEET\_DATA}: The parser is processing the main worksheet data (rows and cells).
%  \item \texttt{AFTER\_SHEET\_DATA}: The parser has detected the closing \texttt{</sheetData>} and is collecting any subsequent chunks until the end of the file.
%  \item \texttt{END\_OF\_FILE}: The parser has reached the end of the stream, and no further parsing is needed.
%\end{itemize}
%
%\paragraph{1. In between nodes}
%In \texttt{BEFORE\_COLS}, \texttt{BEFORE\_SHEET\_DATA}, or \texttt{AFTER\_SHEET\_DATA}, the parser looks for tag boundaries (e.g., \verb|<cols>|, \verb|<sheetData>|) or the end of file. When it detects such a tag at position \verb|i| in the buffer:
%\begin{enumerate}
%  \item It \emph{cuts} the buffer at \verb|i|, splitting the buffer into two segments:
%   (1) from the start to the end of the located tag, and (2) the bytes after that tag.
%  \item It returns the first segment as a chunk to add to the output, which the parser outputs or processes for the current state.
%  \item The second segment becomes the updated buffer, ensuring the XML node boundary is never split across segments.
%\end{enumerate}
%If the parser fails to find any relevant tag or end of file, it returns no new output chunk, remains in the same state and leaves the buffer unmodified. This ensures that subsequent incoming data can be appended, eventually providing enough content to match the target boundary.
%
%\paragraph{2. In column definition node.}
%In \texttt{IN\_COLS}, the parser searches for \verb|</cols>| or any \verb|<col .../>| elements. On detecting \verb|</cols>|, it cuts the buffer right before the closing tag and transitions to \texttt{BEFORE\_SHEET\_DATA}. If it does not find the closing tag but locates complete \verb|<col .../>| elements, those are consumed and removed from the buffer as a clean cut. If no match is found, the buffer remains unchanged, waiting for more data to complete the matching.
%
%\paragraph{3.  In sheet data node.}
%While \texttt{IN\_SHEET\_DATA}, the parser similarly looks for \verb|</sheetData>| or row boundaries \texttt{<row ...>...</row>} (or \texttt{<row .../>}). If \verb|</sheetData>| is found, the parser cuts the buffer at that tag and transitions to \texttt{AFTER\_SHEET\_DATA}. If no boundary is present yet, the buffer remains unaltered to combine with future data chunks.
%
%\paragraph{Ensuring ``Clean Cuts'' within the Streaming Buffer.}
%Our parser reads raw XML data into a working buffer and ensures each XML node boundary is cleanly preserved, so no node is split across buffers. Each state-transition function adheres to this principle, only cutting when a relevant XML tag is fully identified. 
%
%\paragraph{Buffer Updates and State Progress.}
%By constraining parsing decisions to these well-defined cuts, the parser never splits an XML node across multiple buffers. Each function (1) cuts on the matching tag boundary and updates the parser state, or (2) leaves the buffer intact if no complete boundary is found, thus allowing the next chunk to complete it. This way the parser can process the XML incrementally without needing to store the entire document in memory and it is ensured that the regular expressions dont't fail due to incomplete XML nodes and also only the regular expressions that are relevant to the current section of the document are applied, incresing the performance in comparison to applying all regular expressions on the whole document.

% eingefügt


\subsection{Preprocessing with Context-Free Pruning}

  


% eingefügt ende

% eingefügt zwei

\subsection{Context-Aware Post-Processing and Format Override Detection}


After pre-processing, the document is significantly reduced in size and complexity:
only those rows and cells remain which could not be safely discarded
based on context-free visual criteria.
These include cells that do not clearly lack visual styling
and rows that are not part of a redundant sequence of hidden entries.

 
 During pre-processing, we removed cells only if they had no background colour or border.
 However, some transparent-looking cells must not be removed,
 because they may intentionally override an inherited fill or border.
 This is often the case when a background is explicitly applied
 to suppress a coloured fill defined at the row or column level.
 Technically, these cells have style flags such as applyFill="1"
 or applyBorder="1" set, which indicates that their visual appearance
 should take precedence over inherited formatting.

Detecting these overrides requires a deeper understanding of Excel's styling cascade
and cannot be handled reliably using simple regular expressions.
Consequently, this logic is deferred to the post-processing phase,
where the document context is fully available
and a cell's styling can be interpreted in relation
to its parent row or column formatting.

The regex to match those cells and capture the relevant attributes
is shown in Listing \ref{lst:regex_rows_with_cells}.
The regex captures the row id, hidden state, style, and height attributes of the row.
The regex also captures the content of the row if it contains any cells.
This allows us to identify rows that are empty or contain only whitespace.

\begin{lstlisting}[caption={Regex for rows in cell range},label={lst:regex_rows_with_cells}]
<row\b
# capture the whole attribute string
( # capture the row's
  (?:[^>]*\br="([^"]+?)")?  # row id,
  (?:[^>]*\bhidden="(1)")?  # hidden state,
  (?:[^>]*\bs="([^"]+?)")?  # style,
  (?:[^>]*\bht="([^"]+?)")? # and height
  [^/>]*?
)
# match empty rows and rows with cells
(?:/>|>(.*?)</row>)
\end{lstlisting}


%\begin{algorithm}[ht]
%  \caption{Row Replacement Procedure}
%  \label{alg:replace_row}
%  \begin{algorithmic}[1]
%    \STATE $was\_already\_empty \gets (row\_content \text{ is empty})$
%    \IF{not $was\_already\_empty$}
%      \STATE $row\_content \gets CleanCells(row\_content)$
%    \ENDIF
%    \STATE $is\_empty \gets (row\_content \text{ is empty})$
%    \STATE $became\_empty\_by\_cleaning \gets (\lnot was\_already\_empty) \land is\_empty$
%    
%    \IF{$is\_empty$}
%      \IF{$is\_hidden$}
%        \IF{$PreviousRowWasAlsoHidden()$}
%          \RETURN $DeleteRow()$
%        \ELSE
%          \RETURN $ReplaceWithSelfClosingRow()$
%        \ENDIF
%      \ENDIF
%      \IF{$\left((\lnot has\_style) \lor (\lnot has\_fill\_or\_border)\right) \land \lnot has\_height$}
%        \RETURN $DeleteRow()$
%      \ENDIF
%      \IF{$became\_empty\_by\_cleaning \land has\_height$}
%        \RETURN $AddCustomHeightIfAbsent()$
%      \ENDIF
%      \RETURN $ReplaceWithSelfClosingRow()$
%    \ENDIF
%
%    \RETURN $DoNotReplaceRow()$
%  \end{algorithmic}
%\end{algorithm}



% eingefügt zwei ende


% eingefügt drei

\subsection{Two-Tiered Post-Processing for Efficient Truncation}


To further optimise the document reduction process,
the post-processing phase is internally divided into two distinct stages.
This separation is motivated by the observation
that identifying rows containing cells is computationally more expensive
than identifying empty rows.
Specifically, matching the full structure of <row> elements
that contain <c> (cell) definitions requires more complex regular expressions
compared to detecting empty rows, which can be handled more efficiently.

In the first stage of post-processing,
we focus exclusively on rows that contain cell definitions.
By scanning for the final occurrence of a <c> element
and then locating the surrounding <row> end tag,
we determine the upper boundary of the last row containing meaningful cell data.
The document is truncated at this point, ensuring that no superfluous cell-containing rows are preserved.

In the second stage, we analyse rows that do not contain any cells.
Here, we identify the last row that is visually relevant based on lightweight heuristics:
the presence of a visible background colour or border,
or the absence of the hidden="1" attribute.
These attributes can be detected efficiently through simple pattern matching.
Once this second visual boundary is found, the document is truncated again.
All subsequent rows are excluded from further analysis,
significantly reducing the processing time,
as expensive regex-based matching is no longer necessary on these parts of the document.

By separating the treatment of cell-containing and cell-less rows,
we are able to apply specialised heuristics and optimise pattern matching logic,
resulting in both performance improvements
and robust preservation of visually meaningful content.

\subsection{Rule-Based Removal of Empty Rows with Structural and Visual Constraints}

In the post-processing step focused on rows and cells containing content,
we apply a rule-based filtering mechanism implemented using regular expressions.
The process begins by iterating over all row elements (<row>)
in the XML representation of the worksheet.
Each row is evaluated against a set of structural and visual criteria
to determine whether it can be safely removed.

A row is only eligible for removal if all of the following conditions are met:

It contains no cells (i.e. no <c> elements present).

It is not marked as hidden, or if it is hidden, the immediately preceding row must not be hidden. If the preceding row is hidden however, the row after that can safely be removed if it also has no content. This condition ensures that we do not break intended collapsible groups of visible data.

It does not contain any visible styling such as background colour or borders.
The presence of such styling suggests that the row contributes to the visual appearance of the sheet and must be preserved.

If the row contains cells, these cells must first be examined if they can be removed,
if so, the row can be deleted as well, otherwise, the row will be kept

\subsection{Cell-Level Evaluation and Visual Override Resolution}

Before a row can be safely removed, its constituent cells must be individually examined. As part of the initial regular expression match that identifies candidate <row> elements, the styling information of the row itself is extracted. This context is critical when evaluating the visual significance of each cell within the row.

Each cell (<c>) is analysed according to the following decision logic:

Content Check: If the cell contains non-empty content, it is retained immediately.

Whitespace Detection: If the cell contains a string, we check whether it consists solely of whitespace characters. This includes identifying that the cell uses a shared string by matching the t="s" attribute and if so resolving shared strings via looking up the index inside the <v> element which is the child of the <c> element and lookup which string in the SharedStrings table it references. If the corresponding string is pure whitespace, the cell content is cleared and the cell is flagged as empty.

Styling Context: Once a cell is deemed empty, we assess whether it contributes visually through styling:

If the cell defines a background colour or border, it is preserved.

If not, we then check whether it overrides inherited styling from its parent row or column. A white background or explicit border in the cell may overwrite a coloured row or column style. In such cases, even visually “empty” cells must not be removed, as doing so would alter the rendered appearance of the worksheet.

Safe Deletion: Only if the cell neither contains meaningful content, nor introduces visual styling, nor overrides inherited formats, can it be safely removed.

This process ensures that cells are not only semantically empty, but also visually redundant before they are discarded. By incorporating row-level styling context captured during the earlier regex match and column-level styling that we extracted earlier in the <cols> node, we avoid unnecessary recomputation and ensure correctness across nested formatting layers.
% eingefügt drei ende




\paragraph{Putting It All Together.}
Through these transitions, the parser identifies the \texttt{<cols>} segment (if any) and the main \texttt{<sheetData>} block (if any). After discovering \texttt{</sheetData>}, it enters a final phase to read any remaining data. Upon reaching the end of the stream, it moves to \texttt{END\_OF\_FILE}. Then it does postprocessing to the collected output chunks that each state outputted in a different output buffer, so that the postprocessing can be more efficient by only applying the post processing that targets specific regions of the files only get applied to that paticular file section and not the whole document.  Later sections detail how, \emph{within} each state, cells, rows, and styles are selectively removed or retained to produce a more compact representation.

\paragraph{Data Cleaning within the \texttt{sheetData} Node.}
Most of the significant size reduction occurs in the \texttt{sheetData} node, which often contains the primary causes of file bloat. Our data cleaning proceeds as follows:

\noindent
\textbf{Overview.} We streamline the XML-based spreadsheet by removing extraneous cell entries and entire rows that do not contribute to visible content or formatting. In particular, we remove cells containing only whitespace, those referencing styles that provide only a basic font (i.e., no genuine fill or border), and rows with no meaningful content.

\subsection{Automatic Deletion of White Spaces}

Excel’s Open XML format manages textual data through a shared string table stored in \texttt{sharedStrings.xml}. Identical strings appearing multiple times are stored only once, but whitespace-only entries can also reside in this table and cells referencing those whitespaces  bloat the file.  Such a shared strings table would loke like as shown Listing~\ref{lst:sharedStringsSst}. If a cell is declared in the xml like this \lstinline|<c r="A1" t="s"><v>42</v></c>| it means that the attribute t="s" specifies that it is a shared string and it holds the index of the referenced shared string inside the <v> element which is index 42 here. Inside the shared strings File that shared string is the 43th <si> element inside the <sst> node because the index starts at 0. That cell would show "Hello VLDB" to the user. However a cell referencing the shared string with the index 43 would contain just a whitespace character, because the content of the \texttt{<t>} element is empty, indicating that no text is stored. The attribute \texttt{xml:space="preserve"} suggests that this represents a whitespace character. Therefore, all cells referencing index \texttt{43} for the reusable string can be deleted, as they contain no visible content.

\begin{lstlisting}[caption={Shared String Table (\texttt{sst}) in \texttt{sharedStrings.xml}}, label={lst:sharedStringsSst}, firstnumber=48]
<sst count="17962" uniqueCount="732">
  <!-- 42 other si elements  -->
  <si>
    <t>Hello ICDE</t>
  </si>
  <si>
    <t xml:space="preserve"></t>
  </si>
  <!-- ... -->
</sst>  
\end{lstlisting}



\subsubsection{Pattern-Based Detection and Removal of Whitespaces}
We first scan the shared string table to detect entries consisting solely of whitespace. Indices of these redundant entries are marked for later removal. We compile these indices into a single regular expression pattern, matching \texttt{<v>} elements in the worksheet XML that reference them. All such occurrences are removed, thus clearing the affected cell entries.

%\begin{verbatim}
%'<v>(?:43|47|53)</v>'
%\end{verbatim}

%\subsubsection{Application to Sheet Data}
%This regular expression is applied to the raw XML to locate and remove all \texttt{<v>} elements referencing whitespace-only strings. The removal cuts file size significantly without modifying meaningful content. And it also simplifies regular expressions that will be applied next, because now all whitespace cells can be matched easily because they dont contain a child <v> node anymore and can hence be matched with a regular expression that is both less complex and thus faster and safer because there is less likely to overmatch.



%\section{Elimination of Invisible Cells and Removal of Redundant XML Elements}

%Cells in Excel's Open XML reference styles (defined in \lstinline|styles.xml| under \lstinline|cellXfs|) may sometimes only set font properties---without explicit fill or border attributes---and therefore can be considered unnecessary. Removing these cells reduces file size and simplifies further processing.

%\subsection{Elimination of Invisible Cells}

%\subsubsection{Identifying Minimal Style Definitions}

%We parse the \lstinline|cellXfs| list to locate style definitions that lack custom fill or border settings. In other words, these definitions have:
%\begin{itemize}
%  \item \lstinline|fill_id| and \lstinline|border_id| either unset or set to zero.
%  \item \lstinline|apply_fill| and \lstinline|apply_border| flagged as \lstinline|False|.
%\end{itemize}

%\subsubsection{Regular Expression Construction for Cell Matching}

%Once we have the indices of minimal styles, we build a regular expression that detects cell elements (\lstinline|<c>|) whose style attribute (\lstinline|s|) corresponds to one of these indices.

%\paragraph{Regex Listing for Minimal Cell Matching}
%\begin{lstlisting}[caption={Regex for Minimal Cell Matching}, label={lst:regex-minimal}]
%<c\b(?=[^>]*\bs="(?:7|11|17)")% 
%[^>]*(?:\/>|>\s*<\/c>)
%\end{lstlisting}

% \paragraph{Pattern Breakdown:}
% \begin{itemize}
%   \item \textbf{1:} \lstinline|<c\b| \\
%     Matches the opening of a cell element (\lstinline|<c>|). The word boundary (\lstinline|\bs|) ensures that similar tag names are not mistakenly matched.
%   \item \textbf{2:} \lstinline|(?=[^>]*\bs="(?:7|11|17)")| \\
%     A positive lookahead that checks for the presence of a style attribute (\lstinline|s|) whose value is one of the minimal style indices (here, \lstinline|7|, \lstinline|11|, or \lstinline|17|). It does not match if the attribute is missing or does not match one of these values.
%   \item \textbf{3:} \lstinline{[^>]*(?:\/>|>\s*<\/c>)} \\
%     Matches the remainder of the cell tag, accommodating both self-closing tags (\lstinline|/>|) and tags with explicit closing elements (\lstinline|</c>|). This component excludes malformed tags or tags with unexpected nested content.
% \end{itemize}

% \subsubsection{Substitution to Remove Redundant References}

% Applying the above pattern to the XML buffer, all cell elements referencing these minimal styles are removed. This includes cells that are already empty (lacking a nested \lstinline|<v>| node) or cells whose content has been removed in a previous processing step. Removing these cells reduces file size and simplifies subsequent processing, as rows with no remaining cells can be eliminated with an additional regular expression.

%\subsection{Deleting Invisible Cells that Depend on Row or Column Styles}

%Even after the initial clean-up, some cells might persist with no truly visible style. Since Excel's style system can inherit properties from row or column settings, we must compare each cell's local style properties (e.g., \lstinline|apply_fill|, \lstinline|apply_border|) with those inherited from its row or column.

%\paragraph{Checks Include:}
%\begin{itemize}
%  \item \textbf{Border:} If a cell has \lstinline|apply_border| set, its \lstinline|border_id| is compared against the row and column \lstinline|border_id|. If they match and the border object is empty, the cell is deemed invisible with respect to borders.
%  \item \textbf{Fill:} If \lstinline|apply_fill| is set, the cell's \lstinline|fill_id| is compared with that of its row or column. A match, or a fill object with a pattern type of \lstinline|none|, implies that no unique fill is applied.
%\end{itemize}
%Any cell with no distinct fill or border (and without any value) is removed.

% \subsubsection{Regular Expression for Identifying Invisible Cells Dependent on Row/Column Styles}

%\paragraph{Regex Listing for Invisible Cell Detection}
% \begin{lstlisting}[caption={Regex for Invisible Cell Detection}, label={lst:regex-invisible}]
% <c\b(?=[^>]*\br="([A-Za-z]+)(\d+)")% 
% (?=[^>]*\bs="(?:3|5)")% 
% [^>]*(?:\/>|>\s*<\/c>)
% \end{lstlisting}

%\paragraph{Pattern Breakdown:}
%\begin{itemize}
%  \item \textbf{1:} \lstinline|<c\b| \\
%    Matches the beginning of a cell element.
%  \item \textbf{2:} \lstinline|(?=[^>]*\br="([A-Za-z]+)(\d+)")| \\
%    A positive lookahead that captures the cell reference attribute (\lstinline|r|), splitting it into its alphabetical and numeric parts (e.g., \lstinline|A1| becomes \lstinline|A| and \lstinline|1|). It does not match if the cell reference is missing or improperly formatted.
%  \item \textbf{3:} \lstinline|(?=[^>]*\bs="(?:3|5)")| \\
%    A positive lookahead that checks for a style attribute (\lstinline|s|) matching the criteria for inherited fill or border properties (here, the allowed values are \lstinline|3| or \lstinline|5|). If the style does not match, the cell is not considered invisible by this rule.
%  \item \textbf{4:} \lstinline{[^>]*(?:\/>|>\s*<\/c>)} \\
%    Matches the remainder of the cell tag, accommodating both self-closing and paired tags.
%\end{itemize}

% \subsection{Removal of Empty Rows in Worksheet XML}

% After removing invisible cells, rows that no longer contain any cell elements can be eliminated. These empty \lstinline|<row>| elements, which may lack fill, border, or even height attributes, do not contribute to displayed content but can unnecessarily increase file size.

% \subsubsection{Regular Expression for Identifying Empty Rows}

% \paragraph{Regex Listing for Empty Row Removal}
% \begin{lstlisting}[caption={Regex for Empty Row Removal}, label={lst:regex-empty}]
% <row(?![^>]*\bht=)%
% (?=[^>]*(?:\bs="
% (?:7|11|17)
% "|(?:(?!\bs=").)*))%
% [^>]*?(?:\/>|></row>)
% \end{lstlisting}

%\paragraph{Pattern Breakdown:}
%\begin{itemize}
%  \item \textbf{1:} \lstinline|<row| \\
%    Matches the opening of a \lstinline|<row>| tag.
%  \item \textbf{2:} \lstinline|(?![^>]*\bht=)| \\
%    A negative lookahead ensuring that no height attribute (\lstinline|ht|) is present. Rows with a height attribute are excluded, as their height may affect layout.
%  \item \textbf{3:} \lstinline{(?=[^>]*(?:\bs="} \\
%    \item \textbf{4:} \lstinline{(?:7|11|17)} \\
%    \item \textbf{5:} \lstinline{"|(?:(?!\bs=").)*))} \\
%  \item \textbf{6:} \lstinline|[^>]*?| \\
%    A minimal match for the remainder of the tag.
%  \item \textbf{7:} \lstinline|(?:\/>|></row>)| \\
%    Matches the end of the row element, handling both self-closing tags (\lstinline|/>|) and paired row tags (\lstinline|</row>|).
%\end{itemize}

% \subsection{Bounding Rows and Columns}
% \label{sec:bounding_rows_and_columns}
% After purging cells without visible content or style, only cells with meaningful data or formatting remain. We then define the boundaries for rows and columns:
%\begin{itemize}
%  \item \textbf{Rows} below the last cell that contained data or had a style that is visible to the user are removed. Even if a row has a height attribute, it does not affect the document's appearance if there is no content below.
%  \item \textbf{Columns} to the right of the last cell that contained data or had a style that is visible to the user a defined maximum (e.g., \lstinline|MaxCol|) are similarly eliminated from the \lstinline|<cols>| section.
%\end{itemize}
%For example, if the last row with visible content is row 42, the regular expression is constructed to match any row identifiers indicating rows beyond this threshold. The regex is designed to match IDs that are longer than three characters (i.e., row numbers from 100 upwards) as well as rows in the range 43--99.

% \subsubsection{Regular Expression for Bounding Rows}

% \paragraph{Regex Listing for Row Boundary Matching}
% \begin{lstlisting}[caption={Regex for Row Boundary Matching}, label={lst:regex-boundary}]
% <row[^r]+r="
%(:\d{3,}|4[3-9]|[5]\d{1})
%"[^/>]*?(?:\/>|></row>)
%\end{lstlisting}

%\paragraph{Pattern Breakdown:}
%\begin{itemize}
%  \item \textbf{1:} \lstinline|<row[^r]+r="| \\
%    Matches a \lstinline|<row>| tag that contains an attribute \lstinline|r| indicating the row number.
%  \item \textbf{2:} \lstinline|(:\d{3,}|4[3-9]|4\d{1})| \\
%    Captures row numbers that fall outside the defined boundary:
%    \begin{itemize}
%      \item \lstinline|\d{3,}| matches any row number with three or more digits (i.e., 100 or above). When used inline, the curly braces must be escaped, so it becomes \lstinline|\d\{\3,\}|.
%      \item \lstinline|4[3-9]| matches rows numbered 43 through 49.
%      \item \lstinline|4\d{1}| captures a  4 and as many additional digit as are inside the \d{} quantifier, ins this case just one more 1 digit.  
%    \end{itemize}
%  \item \textbf{3:} \lstinline|[^/>]*?(?:\/>|></row>)| \\
%    Matches the remainder of the row element, accommodating both self-closing tags (\lstinline|/>|) and paired tags (\lstinline|</row>|).
%\end{itemize}

%Similarly if the max row index would be 123 group would like like this: \lstinline|(:\d{4,} || \lstinline|[2-9]\d{2} || \lstinline|1[3-9]\d{1} || \lstinline|12[4-9])| 
% and would match any row number that is 114 or above.


  


% \section{Methodology}
%
% \subsection{Initial Experiments with XML Parsers}
% 
% \subsubsection{DOM-Based Prototype}
% Our first prototype employed a DOM parser to read and manipulate the XML subfiles of an Excel workbook (essentially a ZIP containing multiple XML documents).
% 
% \paragraph{Rationale.}
% A DOM-based approach is straightforward to implement and allows direct in-memory manipulation of the entire document tree.
% 
% \paragraph{Key Observation.}
% For very large Excel files (unzipped size in the hundreds of MB or more), this approach heavily taxed both runtime and memory. A single thread could consume all available RAM, making multithreading infeasible. Despite these drawbacks, the DOM-based prototype was a successful \textit{proof of concept} that demonstrated the general feasibility of performing the necessary bloat cleanups on Excel files’ XML data.
% 
% \subsubsection{Switching to Streaming Parsers}
% To reduce memory usage, we tested a series of \emph{streaming} or event-driven XML parsers (e.g., StAX, SAX, or custom “InterParse”-like approaches):
% 
% \begin{itemize}
%     \item \textbf{Outcome:} These approaches offered modest speed improvements and somewhat lower memory footprints than DOM, but still proved insufficient for very large documents—especially in Python—because partial trees were still built in memory.
%     \item \textbf{Further Variation:} Even minimal-memory parsers with event-based processing used little RAM but introduced large numbers of parse events, reducing overall performance.
% \end{itemize}
% 
% At this stage, it became clear that strictly hierarchical parsers might always be at a disadvantage for large-scale cleanup, suggesting we test whether XML’s structure truly requires a tree-based model or if a more direct, regex-based solution could suffice.
% 
% \subsection{Proof-of-Concept Using Regular Expressions}
% 
% \subsubsection{Processing the Byte Stream}
% We developed a \textbf{proof-of-concept} (PoC) that treats each XML subfile as a raw byte stream, applying \textit{regular expressions} without fully decoding or storing the entire document in memory. The approach processes:
% 
% \begin{enumerate}
%     \item \emph{Inner Nodes First} (e.g., cells) to remove unnecessary or oversized elements,
%     \item \emph{Outer Nodes Next} (e.g., entire rows) in a subsequent pass once the inner nodes are removed.
% \end{enumerate}
% 
% By iterating “from the inside out,” we flatten or remove large parts of the hierarchy. Comparing the output to the DOM-based solution confirmed that both removed identical bloat elements.
% 
% \subsubsection{Chunk-Based Parsing}
% To handle large XML files without loading them in entirety:
% 
% \begin{itemize}
%     \item \textbf{Chunking:} We split each XML file into chunks of manageable size, ensuring boundaries align with valid XML tags.
%     \item \textbf{Benefit:} Memory usage drops because we do not keep the entire uncompressed XML in RAM.
%     \item \textbf{Challenge:} Chunk boundaries can complicate regex usage if tags are split; small buffers or additional passes mitigate this issue.
% \end{itemize}
% 
% \subsubsection{Multi-Pass vs.\ Single-Pass Considerations}
% Because certain deletions require global context (e.g., row height that only matters if subsequent rows exist), our initial PoC used multiple consecutive regex passes:
% 
% \begin{itemize}
%     \item \emph{Pass 1:} Remove easily detected bloat such as empty cells or irrelevant attributes.
%     \item \emph{Pass 2:} Remove rows that are definitively empty once their cells are gone.
%     \item \emph{Pass 3:} Adjust elements requiring knowledge of the entire sheet, such as row or style attributes.
% \end{itemize}
% 
% Though this worked well for most files, it caused significant \emph{runtime outliers} in some cases. Careful \emph{regex optimizations} (avoiding catastrophic backtracking) drastically reduced these outliers.
% 
% \subsection{Optimizing the Regex Approach}
% 
% \subsubsection{Avoiding Catastrophic Backtracking}
% Experiments revealed that seemingly small regex details can lead to huge performance impacts:
% 
% \begin{itemize}
%     \item Multiple \lstinline|.*?| segments or complex capturing groups often triggered heavy backtracking.
%     \item Sometimes a \emph{negative character class} like \lstinline|[^"]+| was faster than \lstinline|\d+| or \lstinline|[A-Za-z]+|, depending on the matching engine’s internal logic.
%     \item Parsing complex attributes separately (first capturing the entire attribute string, then dissecting it) could outperform capturing them directly with multiple lookaheads.
% \end{itemize}
% 
% Ultimately, carefully minimizing and structuring the regex components, with limited use of \lstinline|.*?|, proved most efficient in common scenarios.
% 
% \subsubsection{Nested vs.\ Global Scanning}
% We compared:
% 
% \begin{enumerate}
%     \item \textbf{Global:} A single pass to match and remove all relevant tags (e.g.\ rows, cells) at once.
%     \item \textbf{Nested:} Process each row first, then parse its cells within a smaller substring.
% \end{enumerate}
% 
% Although global matching was initially expected to reduce overhead, \emph{nested} matching proved \textbf{faster} and more elegant. It allowed row-level information (e.g., style or height) to be used immediately when deciding which cells to remove, reducing additional passes and large-scale backtracking.
% 
% \subsection{Multithreading and Task Distribution}
% 
% \subsubsection{Early Proof-of-Concept}
% An initial design spawned one thread per Excel file, each of which also launched one thread per worksheet. This left many CPU cores idle when the remaining tasks were fewer than the available threads.
% 
% \subsubsection{Final Multithreading Architecture}
% In the improved approach, we create a fixed number of threads (one per CPU core) and divide the workflow into two phases:
% 
% \begin{enumerate}
%     \item \textbf{Analysis Phase:} Each thread scans assigned files for workbook metadata, generating tasks for each worksheet or other sub-components.
%     \item \textbf{Execution Phase:} These tasks are processed in parallel, irrespective of which file they belong to, ensuring no core remains idle once it finishes a task.
% \end{enumerate}
% 
% Once a sheet is finished, a coordinating thread writes it into the Excel ZIP container. This can occur in parallel for multiple workbooks, but typically concurrency is balanced to avoid excessive context switching.
% 
% \subsection{Summary of the Methodological Approach}
% \begin{enumerate}
%     \item \textbf{DOM-Based Parsing} demonstrated feasibility but was too memory-intensive for large files.
%     \item \textbf{Streaming Parsers} offered some improvements but not enough for truly massive workbooks.
%     \item \textbf{Regex Proof-of-Concept} confirmed that XML can be processed \emph{without} fully building a hierarchical model.
%     \item \textbf{Chunk-Based + Multi-Pass} handling reduced memory usage and managed complex global dependencies.
%     \item \textbf{Regex Fine-Tuning} eliminated costly backtracking, sometimes preferring nested over global scanning.
%     \item \textbf{Multithreading Redesign} used a fixed-thread-pool to ensure full CPU utilization and minimized idle cores.
% \end{enumerate}
% 
% Overall, this multi-stage approach proved both scalable and performant for “bloat removal,” offering a viable alternative to traditional XML tree-based methods for large Excel files while keeping memory overhead manageable.
% 
%
%
%\subsection{Reporting Reduction Results}
%
%To provide a breakdown of how many bytes were saved by deleting whitespaces, empty cells, columns, and rows, we compare the file size before and after each reduction step. For this purpose, we used Excel Schrink Analyser to calculate the size of the worksheet by converting the XML structure into a string, as it would be stored in the file, and compute the length of this string. We then perform a reduction step and again calculate the length of the new string to determine the difference caused by each reduction. Finally, we calculate the percentage reduction relative to the original size. These results are saved in a CSV table.~\cite{excel_shrink_results_csv}
%


 
